\documentclass[a4paper,12pt]{paper}

\usepackage{mycommands}

% Worked-example layout: boxed problem statement followed by a solution.
\newcounter{example}
\renewenvironment{example}[1]{%
  \refstepcounter{example}%
  \begin{mdframed}[linewidth=1pt,innertopmargin=6pt,innerbottommargin=6pt]%
  \noindent\textbf{Example \theexample: #1.}\ }{%
  \end{mdframed}}
\newenvironment{solution}{\par\noindent\textbf{Solution.}\ }{\par\medskip}

\title{Worked examples for Lecture 20: Self-gravitation and rotation}
\author{David Al-Attar \\ Michaelmas Term, \the\year}

\begin{document}

\maketitle

\noindent These examples accompany Lecture~20. They are intended as practice in the concrete
calculations that underlie the lecture, and are less demanding than the questions on the
problem sets. Numerical results were checked with the scripts in the \texttt{python}
directory. Throughout we take $G=6.674\times10^{-11}\,\mathrm{m^{3}\,kg^{-1}\,s^{-2}}$, and for the
Earth a mass $M=5.97\times10^{24}\,\mathrm{kg}$, a mean radius $b=6371\,\mathrm{km}$ and a
rotation rate $\Omega=7.29\times10^{-5}\,\mathrm{s^{-1}}$.

\begin{example}{The homogeneous sphere}
A sphere of radius $b$ has uniform density $\rho$ and total mass $M=\frac{4}{3}\pi\rho b^{3}$.
Solve Poisson's equation to find the gravitational potential $\phi$ and acceleration $\mathbf{g}$
inside and outside the sphere. Compute the gravitational binding energy $\mathcal{V}_{g}$ in two
ways, using both expressions obtained in the lecture, and show that
$\mathcal{V}_{g}=-\frac{3}{5}GM^{2}/b$. Evaluate the result for the Earth and compare it with the
energy released by the largest earthquakes.
\end{example}

\begin{solution}
We take the reference body to be the sphere itself, so that the motion is the identity, the spatial
and referential densities coincide, and $\zeta=\phi$, $\bm{\gamma}=\mathbf{g}$. By symmetry $\phi$
depends only on $r=\|\mathbf{x}\|$, and Poisson's equation $\nabla^{2}\phi=4\pi G\varrho$ reads
\begin{equation}
\frac{1}{r^{2}}\frac{\ddns}{\ddns r}\left(r^{2}\frac{\ddns\phi}{\ddns r}\right)=
\begin{cases}
4\pi G\rho, & r<b,\\
0, & r>b.
\end{cases}
\label{eq:1}
\end{equation}
Integrating twice, the solution that is regular at the origin and tends to zero at infinity is
\begin{equation}
\phi(r)=
\begin{cases}
\frac{2}{3}\pi G\rho\,r^{2}+A, & r<b,\\
-\frac{C}{r}, & r>b,
\end{cases}
\label{eq:2}
\end{equation}
with $A$ and $C$ constants. Since there are no surface masses, the lecture's continuity conditions
require both $\phi$ and $g_{r}=-\frac{\ddns\phi}{\ddns r}$ to be continuous at $r=b$. Continuity of
the radial derivative gives $\frac{4}{3}\pi G\rho b=C/b^{2}$, hence $C=\frac{4}{3}\pi G\rho b^{3}=GM$,
while continuity of $\phi$ then gives $A=-GM/b-\frac{2}{3}\pi G\rho b^{2}=-GM/b-GM/(2b)$, where we have
used $\rho=3M/(4\pi b^{3})$. Collecting terms,
\begin{equation}
\phi(r)=
\begin{cases}
-\dfrac{GM}{2b^{3}}\left(3b^{2}-r^{2}\right), & r\le b,\\[8pt]
-\dfrac{GM}{r}, & r\ge b,
\end{cases}
\qquad
g_{r}(r)=
\begin{cases}
-\dfrac{GM}{b^{3}}\,r=-\dfrac{4}{3}\pi G\rho\,r, & r\le b,\\[8pt]
-\dfrac{GM}{r^{2}}, & r\ge b.
\end{cases}
\label{eq:3}
\end{equation}
The exterior field is that of a point mass $M$ at the origin, while inside the sphere the
acceleration grows linearly with radius, as only the mass interior to $r$ contributes. The potential
at the centre is $-\frac{3}{2}GM/b$, one and a half times its surface value.

For the binding energy we first use the volume form $\mathcal{V}_{g}=\frac{1}{2}\int_{M_{t}}\varrho\phi\dd^{3}\mathbf{y}$
from the lecture, where the integral extends only over the body. With $\dd^{3}\mathbf{y}=4\pi r^{2}\dd r$,
\begin{equation}
\begin{split}
\mathcal{V}_{g}&=\frac{1}{2}\rho\int_{0}^{b}\left[-\frac{GM}{2b^{3}}\left(3b^{2}-r^{2}\right)\right]4\pi r^{2}\dd r\\
&=-\frac{\pi\rho GM}{b^{3}}\left(b^{5}-\frac{b^{5}}{5}\right)=-\frac{4\pi\rho GMb^{2}}{5}=-\frac{3}{5}\frac{GM^{2}}{b},
\end{split}
\label{eq:4}
\end{equation}
where the last step again uses $\rho=3M/(4\pi b^{3})$. Next we use the field-energy form
$\mathcal{V}_{g}=-\frac{1}{8\pi G}\int_{\mathbb{R}^{3}}\|\nabla\phi\|^{2}\dd^{3}\mathbf{y}$, for which
the integral extends over all space. Here $\|\nabla\phi\|^{2}=g_{r}^{2}$, and
\begin{equation}
\begin{split}
\int_{\mathbb{R}^{3}}\|\nabla\phi\|^{2}\dd^{3}\mathbf{y}
&=4\pi\left(\frac{G^{2}M^{2}}{b^{6}}\int_{0}^{b}r^{4}\dd r+G^{2}M^{2}\int_{b}^{\infty}\frac{\dd r}{r^{2}}\right)\\
&=4\pi G^{2}M^{2}\left(\frac{1}{5b}+\frac{1}{b}\right)=\frac{24\pi}{5}\frac{G^{2}M^{2}}{b},
\end{split}
\label{eq:5}
\end{equation}
so that
\begin{equation}
\mathcal{V}_{g}=-\frac{1}{8\pi G}\cdot\frac{24\pi}{5}\frac{G^{2}M^{2}}{b}=-\frac{3}{5}\frac{GM^{2}}{b},
\label{eq:6}
\end{equation}
in agreement with eq.~(\ref{eq:4}). It is worth noticing how differently the two formulae distribute
the same total: in eq.~(\ref{eq:5}) five sixths of the field energy resides \emph{outside} the sphere,
whereas eq.~(\ref{eq:4}) has no exterior contribution at all. The two expressions differ by a
divergence, and only their integrals agree.

For the Earth, with the values quoted in the preamble, $GM^{2}/b=3.73\times10^{32}\,\mathrm{J}$ and hence
\begin{equation}
\mathcal{V}_{g}=-\frac{3}{5}\frac{GM^{2}}{b}\approx-2.24\times10^{32}\,\mathrm{J}.
\label{eq:7}
\end{equation}
For comparison, the same data give a mean density $3M/(4\pi b^{3})\approx5510\,\mathrm{kg\,m^{-3}}$ and a
surface gravity $GM/b^{2}\approx9.82\,\mathrm{m\,s^{-2}}$. The energy radiated by the largest recorded
earthquakes, such as the 1960 Chile event, is of order $10^{18}$--$10^{19}\,\mathrm{J}$, so the
gravitational binding energy exceeds the energy of the largest earthquake by a factor of
$10^{13}$--$10^{14}$. Because the real Earth's mass is concentrated towards its centre, the true
binding energy is somewhat larger in magnitude than the homogeneous estimate. The comparison makes
plain why gravity cannot be ignored in the long-period dynamics: even a minute fractional change in
the gravitational energy is enormous on the scale of seismic energies.
\end{solution}

\begin{example}{Referential gravity under a rigid motion}
Let $\bphi$ be a motion of the body, and let $\tilde{\bphi}=\mathbf{a}+\mathbf{Q}\bphi$ be the motion
obtained by composing it with a rigid translation $\mathbf{a}$ and rotation $\mathbf{Q}$ (so that
$Q_{ki}Q_{kj}=\delta_{ij}$). Show directly from the integral expressions of the lecture that the
referential gravitational potential is unchanged, $\tilde{\zeta}=\zeta$, while the referential
gravitational acceleration transforms as $\tilde{\gamma}_{i}=Q_{ij}\gamma_{j}$. Comment on the
significance of these results for the Lagrangian and for the equations of motion in a rotating frame.
\end{example}

\begin{solution}
The referential potential is given in the lecture by
\begin{equation}
\zeta(\mathbf{x},t)=-G\int_{M}\frac{\rho(\mathbf{x}')}{\|\bphi(\mathbf{x},t)-\bphi(\mathbf{x}',t)\|}\dd^{3}\mathbf{x}',
\label{eq:8}
\end{equation}
and depends on the motion only through the separation $\bphi(\mathbf{x},t)-\bphi(\mathbf{x}',t)$
of pairs of material points. Under the rigid motion this separation becomes
\begin{equation}
\tilde{\bphi}(\mathbf{x},t)-\tilde{\bphi}(\mathbf{x}',t)
=\mathbf{a}+\mathbf{Q}\bphi(\mathbf{x},t)-\mathbf{a}-\mathbf{Q}\bphi(\mathbf{x}',t)
=\mathbf{Q}\left[\bphi(\mathbf{x},t)-\bphi(\mathbf{x}',t)\right],
\label{eq:9}
\end{equation}
so the translation cancels identically. Writing $\mathbf{d}=\bphi(\mathbf{x},t)-\bphi(\mathbf{x}',t)$,
orthogonality of $\mathbf{Q}$ gives
\begin{equation}
\|\mathbf{Q}\mathbf{d}\|^{2}=Q_{ki}d_{i}Q_{kj}d_{j}=\delta_{ij}d_{i}d_{j}=\|\mathbf{d}\|^{2},
\label{eq:10}
\end{equation}
and hence every denominator in eq.~(\ref{eq:8}) is unchanged. The referential density $\rho(\mathbf{x}')$
is a property of the reference body and does not depend on the motion at all. It follows at once that
$\tilde{\zeta}(\mathbf{x},t)=\zeta(\mathbf{x},t)$.

The referential acceleration is
\begin{equation}
\gamma_{i}(\mathbf{x},t)=-G\int_{M}\frac{\rho(\mathbf{x}')d_{i}}{\|\mathbf{d}\|^{3}}\dd^{3}\mathbf{x}',
\label{eq:11}
\end{equation}
with $\mathbf{d}$ as above. Under the rigid motion $d_{i}\rightarrow Q_{ij}d_{j}$ in the numerator
while, by eq.~(\ref{eq:10}), the denominator is unchanged. Since $\mathbf{Q}$ does not depend on
$\mathbf{x}'$ it can be taken outside the integral, and so
\begin{equation}
\tilde{\gamma}_{i}(\mathbf{x},t)=-G\int_{M}\frac{\rho(\mathbf{x}')Q_{ij}d_{j}}{\|\mathbf{d}\|^{3}}\dd^{3}\mathbf{x}'
=Q_{ij}\gamma_{j}(\mathbf{x},t),
\label{eq:12}
\end{equation}
as required.

The same conclusions can be reached from the spatial description. If the spatial density of the
original motion is $\varrho(\mathbf{y},t)$, that of the rigidly moved body is
$\tilde{\varrho}(\mathbf{y},t)=\varrho[\mathbf{Q}^{\mathrm{T}}(\mathbf{y}-\mathbf{a}),t]$, and because
the Laplacian is invariant under rigid motions the solution of Poisson's equation is simply carried
along, $\tilde{\phi}(\mathbf{y},t)=\phi[\mathbf{Q}^{\mathrm{T}}(\mathbf{y}-\mathbf{a}),t]$.
Evaluating at $\mathbf{y}=\tilde{\bphi}(\mathbf{x},t)$ gives
$\tilde{\zeta}(\mathbf{x},t)=\phi[\bphi(\mathbf{x},t),t]=\zeta(\mathbf{x},t)$, while the chain rule
applied to $\tilde{g}_{i}=-\partial\tilde{\phi}/\partial y_{i}$ produces the factor $Q_{ij}$.

Two consequences are worth drawing out. First, the gravitational binding energy
$\mathcal{V}_{g}=\frac{1}{2}\int_{M}\rho\zeta\dd^{3}\mathbf{x}$ is unchanged by a rigid motion.
The Lagrangian of the lecture is therefore frame indifferent in the same sense as the elastic strain
energy, which is essential if the equations of motion are to be physically meaningful. Second,
taking $\mathbf{a}=\bm{\Phi}(t)$ and $\mathbf{Q}=\mathbf{R}(t)$ to be the translation and rotation
introduced in the lecture for the co-rotating frame, eq.~(\ref{eq:12}) is precisely the identity
$\gamma_{i}=R_{ij}\overline{\gamma}_{j}$ used there, with $\overline{\bm{\gamma}}$ computed from the
relative motion $\overline{\bphi}$. Gravity, unlike the inertial terms, is oblivious to the rotation
of the frame; it is only the acceleration that acquires Coriolis and centrifugal contributions.
\end{solution}

\begin{example}{A steadily rotating frame}
Let $\mathbf{R}(t)$ be the rotation through angle $\Omega t$ about the $\hat{\mathbf{e}}_{3}$ axis.
Compute the matrix $R_{kj}\frac{\ddns R_{kl}}{\ddns t}$, verify that it is anti-symmetric, and use the
lecture's relation $R_{kj}\frac{\ddns R_{kl}}{\ddns t}=\epsilon_{jml}\Omega_{m}$ to identify the angular
velocity vector $\bm{\Omega}$. Write out the Coriolis and centrifugal accelerations in components.
Show that the centrifugal acceleration $\epsilon_{ijk}\epsilon_{klm}\Omega_{j}\Omega_{l}x_{m}$ is
the gradient of $\psi=-\frac{1}{2}\Omega^{2}(x_{1}^{2}+x_{2}^{2})$, and check this against the
general expression $\psi=\frac{1}{2}(\Omega_{i}\Omega_{j}-\Omega_{k}\Omega_{k}\delta_{ij})x_{i}x_{j}$
valid for an arbitrary rotation vector. Evaluate the centrifugal acceleration at the Earth's equator.
\end{example}

\begin{solution}
Writing $c=\cos\Omega t$ and $s=\sin\Omega t$, the rotation matrix and its time derivative are
\begin{equation}
\mathbf{R}(t)=\begin{pmatrix}c&-s&0\\ s&c&0\\ 0&0&1\end{pmatrix},\qquad
\frac{\ddns\mathbf{R}}{\ddns t}=\Omega\begin{pmatrix}-s&-c&0\\ c&-s&0\\ 0&0&0\end{pmatrix}.
\label{eq:13}
\end{equation}
The quantity $R_{kj}\frac{\ddns R_{kl}}{\ddns t}$ is the $(j,l)$ element of
$\mathbf{R}^{\mathrm{T}}\frac{\ddns\mathbf{R}}{\ddns t}$, and multiplying out,
\begin{equation}
\begin{split}
\mathbf{R}^{\mathrm{T}}\frac{\ddns\mathbf{R}}{\ddns t}
&=\Omega\begin{pmatrix}c&s&0\\ -s&c&0\\ 0&0&1\end{pmatrix}
\begin{pmatrix}-s&-c&0\\ c&-s&0\\ 0&0&0\end{pmatrix}\\
&=\Omega\begin{pmatrix}-cs+sc&-c^{2}-s^{2}&0\\ s^{2}+c^{2}&sc-cs&0\\ 0&0&0\end{pmatrix}
=\begin{pmatrix}0&-\Omega&0\\ \Omega&0&0\\ 0&0&0\end{pmatrix},
\end{split}
\label{eq:14}
\end{equation}
which is manifestly anti-symmetric and, notably, independent of time. To identify $\bm{\Omega}$ we
compare with $\epsilon_{jml}\Omega_{m}$. Taking $(j,l)=(1,2)$ gives
$\epsilon_{1m2}\Omega_{m}=\epsilon_{132}\Omega_{3}=-\Omega_{3}$, which must equal $-\Omega$, while
$(j,l)=(2,3)$ gives $\epsilon_{2m3}\Omega_{m}=\Omega_{1}=0$ and $(j,l)=(3,1)$ gives $\Omega_{2}=0$.
Hence $\bm{\Omega}=\Omega\hat{\mathbf{e}}_{3}$, as expected for a rotation about the third axis. More
generally, contracting the lecture's relation with $\epsilon_{jnl}$ and using
$\epsilon_{jml}\epsilon_{jnl}=2\delta_{mn}$ gives the explicit formula
\begin{equation}
\Omega_{n}=\frac{1}{2}\epsilon_{jnl}R_{kj}\frac{\ddns R_{kl}}{\ddns t},
\label{eq:15}
\end{equation}
which recovers the angular velocity from any rotation history $\mathbf{R}(t)$.

With $\bm{\Omega}=\Omega\hat{\mathbf{e}}_{3}$ the Coriolis acceleration $2\epsilon_{ijk}\Omega_{j}\overline{v}_{k}$
is $2\bm{\Omega}\times\overline{\mathbf{v}}$, with components
\begin{equation}
2\epsilon_{ijk}\Omega_{j}\overline{v}_{k}=2\Omega\left(-\overline{v}_{2},\,\overline{v}_{1},\,0\right),
\label{eq:16}
\end{equation}
so it is perpendicular to both the rotation axis and the relative velocity, and vanishes for motion
parallel to the axis. For the centrifugal term we use the identity
$\epsilon_{ijk}\epsilon_{klm}=\delta_{il}\delta_{jm}-\delta_{im}\delta_{jl}$ to write
\begin{equation}
\epsilon_{ijk}\epsilon_{klm}\Omega_{j}\Omega_{l}x_{m}=\Omega_{i}\Omega_{j}x_{j}-\Omega_{j}\Omega_{j}x_{i}
=\left[\bm{\Omega}\times(\bm{\Omega}\times\mathbf{x})\right]_{i},
\label{eq:17}
\end{equation}
and for $\bm{\Omega}=\Omega\hat{\mathbf{e}}_{3}$ this is
\begin{equation}
\epsilon_{ijk}\epsilon_{klm}\Omega_{j}\Omega_{l}x_{m}=\Omega^{2}\left(0,0,x_{3}\right)-\Omega^{2}\left(x_{1},x_{2},x_{3}\right)
=-\Omega^{2}\left(x_{1},\,x_{2},\,0\right),
\label{eq:18}
\end{equation}
which points towards the rotation axis with magnitude $\Omega^{2}$ times the distance from it.

Now let $\psi=-\frac{1}{2}\Omega^{2}(x_{1}^{2}+x_{2}^{2})$. Its gradient is
$\nabla\psi=-\Omega^{2}(x_{1},x_{2},0)$, which by eq.~(\ref{eq:18}) is exactly the centrifugal
acceleration. Thus the term appearing on the left-hand side of the lecture's equations of motion is
$\rho\,\partial\psi/\partial x_{i}$, and when it is carried to the right-hand side it becomes the
centrifugal body force $-\rho\,\partial\psi/\partial x_{i}$ per unit volume, directed away from the axis;
$\psi$ is the \textbf{centrifugal potential}. To check the general formula, differentiate
$\psi=\frac{1}{2}(\Omega_{i}\Omega_{j}-\Omega_{k}\Omega_{k}\delta_{ij})x_{i}x_{j}$, noting that the
coefficient matrix is symmetric, to find
\begin{equation}
\frac{\partial\psi}{\partial x_{i}}=\left(\Omega_{i}\Omega_{j}-\Omega_{k}\Omega_{k}\delta_{ij}\right)x_{j}
=\Omega_{i}\Omega_{j}x_{j}-\Omega_{k}\Omega_{k}x_{i},
\label{eq:19}
\end{equation}
which is eq.~(\ref{eq:17}). Substituting $\bm{\Omega}=\Omega\hat{\mathbf{e}}_{3}$ into the general
formula gives
\begin{equation}
\psi=\frac{1}{2}\left[\Omega^{2}x_{3}^{2}-\Omega^{2}(x_{1}^{2}+x_{2}^{2}+x_{3}^{2})\right]
=-\frac{1}{2}\Omega^{2}(x_{1}^{2}+x_{2}^{2}),
\label{eq:20}
\end{equation}
in agreement with the special case. Note that
$\nabla^{2}\psi=-2\Omega^{2}\neq0$, so unlike the gravitational potential in free space the centrifugal
potential is not harmonic.

At the equator the distance from the axis is $b$, and the centrifugal acceleration has magnitude
\begin{equation}
\Omega^{2}b=(7.29\times10^{-5})^{2}\times6.371\times10^{6}\,\mathrm{m\,s^{-2}}\approx3.39\times10^{-2}\,\mathrm{m\,s^{-2}},
\label{eq:21}
\end{equation}
which is $\Omega^{2}b^{3}/GM\approx3.4\times10^{-3}$, or about one part in $290$, of the surface gravity found
in Example~1. This small ratio is what permits the Earth's equilibrium figure to be treated as a
slightly flattened sphere in the following lecture, where the centrifugal and gravitational potentials
are combined into a single geopotential.
\end{solution}

\begin{example}{Scaling analysis with numbers}
Using the order-of-magnitude estimates of the lecture, with $\overline{\rho}=5000\,\mathrm{kg\,m^{-3}}$
and $\overline{v}=8000\,\mathrm{m\,s^{-1}}$ as typical density and wave speed, evaluate the ratio of
gravitational to elastic forces, the ratio $\Omega T$ of Coriolis to inertial terms, and the ratio of
centrifugal to inertial terms, for (a) a $1\,\mathrm{Hz}$ body wave with $L\sim10\,\mathrm{km}$; (b) a
$100\,\mathrm{s}$ surface wave with $L\sim400\,\mathrm{km}$; and (c) the gravest spheroidal mode
$_{0}S_{2}$, with period $54\,\mathrm{min}$ and $L\sim b$. Tabulate the results and state where gravity
and rotation matter.
\end{example}

\begin{solution}
From the lecture, the ratios of the various terms in the linearised equations of motion to the
elastic and inertial terms are
\begin{equation}
\frac{\text{gravitational}}{\text{elastic}}=\frac{4\pi G\overline{\rho}L^{2}}{\overline{v}^{2}},\quad
\frac{\text{Coriolis}}{\text{inertial}}=\Omega T,\quad
\frac{\text{centrifugal}}{\text{inertial}}=\frac{2\overline{\rho}\Omega^{2}U}{\overline{\rho}U/T^{2}}=2(\Omega T)^{2},
\label{eq:22}
\end{equation}
where the final ratio follows from the lecture's estimate $2\overline{\rho}\Omega^{2}U$ for the
centrifugal term; the factor of two is an upper bound, attained when $\mathbf{u}$ is perpendicular to
$\bm{\Omega}$, and is of no consequence in an order-of-magnitude argument. It is convenient to note the
combination
\begin{equation}
4\pi G\overline{\rho}=4\pi\times6.674\times10^{-11}\times5000\,\mathrm{s^{-2}}=4.19\times10^{-6}\,\mathrm{s^{-2}},
\label{eq:23}
\end{equation}
so that the gravitational ratio is $(L/L_{g})^{2}$ with
$L_{g}=\overline{v}/\sqrt{4\pi G\overline{\rho}}\approx3900\,\mathrm{km}$, while the rotational ratios
involve $T$ compared with $\Omega^{-1}\approx1.37\times10^{4}\,\mathrm{s}\approx3.8\,\mathrm{h}$.
Substituting the three sets of values gives the following table.
\begin{center}
\small
\begin{tabular}{lccccc}
\hline
Case & $T$ (s) & $L$ (km) & grav./elastic & $\Omega T$ & $2(\Omega T)^{2}$ \\
\hline
(a) body wave & $1$ & $10$ & $6.55\times10^{-6}$ & $7.29\times10^{-5}$ & $1.06\times10^{-8}$ \\
(b) surface wave & $100$ & $400$ & $1.05\times10^{-2}$ & $7.29\times10^{-3}$ & $1.06\times10^{-4}$ \\
(c) $_{0}S_{2}$ & $3240$ & $6371$ & $2.66$ & $0.236$ & $0.112$ \\
\hline
\end{tabular}
\end{center}
For the body wave every correction is utterly negligible: gravity enters at the level of a few parts
per million and rotation at the level of one part in $10^{4}$ or below, so the elastodynamics of the
earlier lectures is exact for all practical purposes. For the $100\,\mathrm{s}$ surface wave the
gravitational correction has grown to about one per cent, as expected from the $L^{2}$ scaling,
and although still small it is comparable to the fractional variations in wave speed that surface-wave
tomography seeks to resolve, so gravity is routinely included in surface-wave calculations at these
periods. Rotation, by contrast, remains below one per cent. For $_{0}S_{2}$ the gravitational ratio is
of order unity -- this is the lecture's estimate of about $3$ -- and gravity is as important as
elasticity; indeed, the restoring force for this mode is largely gravitational. The Coriolis ratio
$\Omega T\approx0.24$ is no longer small, and it is responsible for the observable splitting of the
gravest modes into multiplets. The centrifugal ratio $2(\Omega T)^{2}\approx0.1$ is smaller but not
negligible at this period.

The table illustrates the general conclusion of the lecture. Gravity matters for deformations whose
length scale approaches $L_{g}$ of a few thousand kilometres, which for seismic waves means periods of
several hundred seconds and longer, while rotation matters for time scales that approach
$\Omega^{-1}$, i.e.\ for the longest-period normal modes, and is dominant only for phenomena such as
tides with periods comparable to a day.
\end{solution}

\begin{example}{Linearised gravitational acceleration for rigid displacements}
The lecture gives the first-order perturbation to the referential gravitational acceleration as
\begin{equation}
\gamma_{i}^{1}=-G\int_{M}\rho'\left(\frac{\delta_{ij}}{\|\mathbf{x}-\mathbf{x}'\|^{3}}-
\frac{3(x_{i}-x_{i}')(x_{j}-x_{j}')}{\|\mathbf{x}-\mathbf{x}'\|^{5}}\right)(u_{j}-u_{j}')\dd^{3}\mathbf{x}'.
\label{eq:24}
\end{equation}
Show that a rigid translation $\mathbf{u}=\mathbf{u}_{0}$ gives $\bm{\gamma}^{1}=\mathbf{0}$, and
that an infinitesimal rigid rotation $\mathbf{u}=\bm{\theta}\times\mathbf{x}$, with $\bm{\theta}$ a
constant vector, gives $\bm{\gamma}^{1}=\bm{\theta}\times\bm{\gamma}^{0}$, where $\bm{\gamma}^{0}$
is the equilibrium referential acceleration. Explain why both results are to be expected.
\end{example}

\begin{solution}
For the translation, $u_{j}-u_{j}'=u_{0j}-u_{0j}=0$ for every pair of points, so the integrand in
eq.~(\ref{eq:24}) vanishes identically and $\bm{\gamma}^{1}=\mathbf{0}$. Nothing needs to be assumed
about the density or the shape of the body. Physically, $\gamma_{i}(\mathbf{x},t)=g_{i}[\bphi(\mathbf{x},t),t]$
is the acceleration felt by the material particle labelled $\mathbf{x}$. Translating the whole body
translates its gravitational field along with it, and each particle continues to feel exactly the
field it felt before; only in the spatial description, at fixed $\mathbf{y}$, does the field change.
This is the linearised form of the translational invariance established in Example~2.

For the rotation write $\mathbf{r}=\mathbf{x}-\mathbf{x}'$ and $r=\|\mathbf{r}\|$. Then
\begin{equation}
u_{j}-u_{j}'=\epsilon_{jkl}\theta_{k}(x_{l}-x_{l}')=\epsilon_{jkl}\theta_{k}r_{l}=(\bm{\theta}\times\mathbf{r})_{j},
\label{eq:25}
\end{equation}
which is perpendicular to $\mathbf{r}$. Contracting with the kernel in eq.~(\ref{eq:24}),
\begin{equation}
\left(\frac{\delta_{ij}}{r^{3}}-\frac{3r_{i}r_{j}}{r^{5}}\right)(\bm{\theta}\times\mathbf{r})_{j}
=\frac{(\bm{\theta}\times\mathbf{r})_{i}}{r^{3}}-\frac{3r_{i}\,\mathbf{r}\cdot(\bm{\theta}\times\mathbf{r})}{r^{5}}
=\frac{(\bm{\theta}\times\mathbf{r})_{i}}{r^{3}},
\label{eq:26}
\end{equation}
because $\mathbf{r}\cdot(\bm{\theta}\times\mathbf{r})=0$. The second, traceless part of the kernel
therefore contributes nothing, and
\begin{equation}
\gamma_{i}^{1}=-G\int_{M}\rho'\frac{\epsilon_{ikl}\theta_{k}r_{l}}{r^{3}}\dd^{3}\mathbf{x}'
=\epsilon_{ikl}\theta_{k}\left(-G\int_{M}\rho'\frac{x_{l}-x_{l}'}{\|\mathbf{x}-\mathbf{x}'\|^{3}}\dd^{3}\mathbf{x}'\right)
=\epsilon_{ikl}\theta_{k}\gamma_{l}^{0},
\label{eq:27}
\end{equation}
where the bracketed integral is the equilibrium referential acceleration obtained by setting
$\bphi(\mathbf{x})=\mathbf{x}$ in the lecture's integral formula. That is, $\bm{\gamma}^{1}=\bm{\theta}\times\bm{\gamma}^{0}$.

This too is what Example~2 predicts. A finite rotation $\mathbf{Q}$ transforms the referential
acceleration as $\bm{\gamma}\rightarrow\mathbf{Q}\bm{\gamma}$, and for an infinitesimal rotation
$Q_{ij}=\delta_{ij}+s\,\epsilon_{ikj}\theta_{k}+O(s^{2})$, corresponding to the motion
$\varphi_{i}=x_{i}+s\,\epsilon_{ikj}\theta_{k}x_{j}+O(s^{2})$, so that
$\gamma_{i}=\gamma_{i}^{0}+s\,\epsilon_{ikj}\theta_{k}\gamma_{j}^{0}+O(s^{2})$; the first-order term is exactly
eq.~(\ref{eq:27}). The gravitational field is rigidly rotated with the body, and each particle feels
the rotated field. In the non-rotating case, Example~2 shows that a rigid motion of an equilibrium
configuration is again an equilibrium configuration, so rigid translations and rotations must be
zero-frequency solutions of the linearised equations of motion; the two calculations here are the
gravitational part of that statement, and the remaining terms must conspire in the same way. When the
frame rotates the argument no longer applies as it stands, because the centrifugal term
$\epsilon_{ijk}\epsilon_{klm}\Omega_{j}\Omega_{l}u_{m}$ does not vanish for a rigid translation
transverse to the axis, and the rigid-body modes of a rotating planet require separate treatment.
\end{solution}

\end{document}
